\chapter{Candidemia}

\section*{Definition and Overview}
Candidemia is the presence of \textit{Candida} fungi in the bloodstream. It is a serious, usually hospital-acquired infection that most often affects people who are already severely ill. \textit{Candida} is a yeast that normally lives harmlessly on the skin and in the gastrointestinal tract, but when it enters the blood it can trigger a whole-body inflammatory response and can seed to distant organs, including the eyes, kidneys, brain, and heart valves. Candidemia is the most common form of invasive candidiasis and is a medical emergency that requires urgent antifungal treatment and removal of any source of infection.

\section*{Epidemiology}
Candidemia is one of the most common bloodstream infections in hospitalised patients, and its incidence has risen over recent decades as intensive care has become more invasive and immunosuppressive therapies more widespread. It occurs almost exclusively in people who are hospitalised, especially those in intensive care units, and is now a leading cause of nosocomial bloodstream infection alongside staphylococcal and enterococcal bacteraemia. Risk is highest among patients with intravenous catheters, those receiving broad-spectrum antibiotics, people undergoing abdominal surgery, patients receiving total parenteral nutrition, and people with weakened immune systems from chemotherapy, organ transplantation, or diabetes. It is also more common in premature infants and in older adults. The infection is associated with substantial illness and death, with attributable mortality often reported in the range of 30--40\%.

\section*{Aetiology and Pathophysiology}
\textit{Candida} species are the cause, most commonly \textit{Candida albicans}, but increasingly other species including \textit{Candida glabrata}, \textit{Candida parapsilosis}, \textit{Candida tropicalis}, and the emerging multidrug-resistant \textit{Candida auris}. These yeasts colonise the gut, skin, and mucous membranes of many healthy people. Disease occurs when the normal barriers are breached, most commonly through a central venous catheter, and the yeast enters the bloodstream. Factors that allow this include:
\begin{itemize}
    \item \textbf{Intravenous catheters:} the yeast can form a sticky biofilm on the plastic tubing that protects it from both the immune system and antifungal drugs
    \item \textbf{Damage to the gut lining:} from surgery, chemotherapy, or prolonged illness, allowing gut yeasts to translocate across the bowel wall
    \item \textbf{Broad-spectrum antibiotics:} which kill competing bacteria and allow yeasts to flourish
    \item \textbf{Neutropenia:} a very low neutrophil count after chemotherapy or in marrow failure
    \item \textbf{Immunosuppressive medicines} used after transplantation or for autoimmune disease
    \item \textbf{Prematurity:} an immature immune system and less robust skin and gut barriers
\end{itemize}
Once in the blood, \textit{Candida} multiplies and triggers widespread inflammation, which can progress to septic shock. The organism can also deposit in tissues, causing abscesses in organs, infection of heart valves (endocarditis), or infection within the eye (endophthalmitis).

\section*{Clinical Features}
Candidemia often produces few specific signs and can be clinically indistinguishable from a bacterial bloodstream infection:
\begin{itemize}
    \item Persistent fever or, paradoxically, low body temperature
    \item Chills and rigors
    \item Rapid heart rate and rapid breathing
    \item Low blood pressure, which can progress to septic shock
    \item Confusion, agitation, or drowsiness
    \item Worsening organ function, such as kidney, liver, or lung failure
    \item In some cases, signs of metastatic spread: painful red eye, skin lesions, a new heart murmur, or back pain
\end{itemize}
Because the affected patients are usually already unwell in hospital, the diagnosis is usually made on laboratory testing rather than on symptoms alone, and candidemia should be suspected in any deteriorating patient with persistent fever despite antibacterial therapy.

\section*{Differential Diagnosis}
Because the presentation is nonspecific, candidemia must be distinguished from other causes of sepsis in the hospitalised patient:
\begin{itemize}
    \item Bacterial bloodstream infections, including catheter-related bacteraemia
    \item Other fungal bloodstream infections, such as \textit{Aspergillus} or cryptococcosis, particularly in immunocompromised patients
    \item Viral infections, such as herpesvirus reactivation in transplant patients
    \item Non-infectious causes of fever in hospital, including drug reactions, blood transfusions, venous thrombosis, and postoperative inflammation
\end{itemize}

\section*{When to Seek Medical Care}
Candidemia is not a condition that can be managed at home; it arises during hospital care and requires urgent hospital treatment. Anyone at home who develops signs of sepsis, or who has a relative in hospital who becomes acutely unwell, should seek help immediately.
\textbf{Call 000 (or local emergency services) for:}
\begin{itemize}
    \item Fever with shaking chills and a rapid heart rate
    \item Confusion, severe drowsiness, or difficulty waking
    \item Low blood pressure with light-headedness or fainting
    \item Difficulty breathing or reduced urine output
    \item A painful red eye, new skin spots, or a new heart murmur in a person who has recently been treated for candidemia
\end{itemize}

\section*{Diagnosis and Investigations}
\begin{itemize}
    \item \textbf{Blood cultures:} the definitive test; multiple bottles are taken from different sites and may need to be repeated, since fungal growth is often slow
    \item \textbf{Species identification:} essential to choose the most effective antifungal and to detect resistant species such as \textit{Candida auris}
    \item \textbf{Culture of the catheter tip:} after removal, if an intravenous catheter is the suspected source
    \item \textbf{Eye examination:} an ophthalmologist checks the retina and vitreous for fungal deposits, ideally within a few days of diagnosis
    \item \textbf{Imaging:} CT or ultrasound of the abdomen, and echocardiography of the heart, to look for abscesses or valve infection
    \item \textbf{Other cultures:} urine, wound, or drainage samples where metastatic spread is suspected
    \item \textbf{Antifungal susceptibility testing:} especially for species known to be resistant to azoles
\end{itemize}

\section*{Management}
Candidemia is treated in hospital, usually in an intensive care or high-dependency unit, and involves:
\begin{itemize}
    \item \textbf{Prompt antifungal therapy:} an intravenous antifungal drug, usually an echinocandin such as micafungin, caspofungin, or anidulafungin, started early and empirically where suspicion is high
    \item \textbf{Removal of the intravenous catheter} whenever possible, since catheter-associated biofilms are difficult to eradicate
    \item \textbf{Source control:} drainage of abscesses and removal of any other infected devices
    \item \textbf{Supportive care:} intravenous fluids, oxygen, blood pressure support with vasopressors, and management of organ failure
    \item \textbf{Treatment of the underlying problem:} such as controlling blood sugar or supporting the immune system
    \item \textbf{Step-down to oral antifungal treatment} once the patient improves, with the total course usually lasting about two weeks after the first negative blood culture
    \item \textbf{Follow-up:} repeat blood cultures to confirm clearance and eye checks to detect late retinal infection
\end{itemize}

\section*{Complications}
\begin{itemize}
    \item Septic shock and multi-organ failure
    \item Endocarditis: infection of the heart valves, which may require prolonged treatment and sometimes surgery
    \item Endophthalmitis: fungal infection inside the eye that can threaten sight
    \item Abscesses in the liver, spleen, kidneys, or brain
    \item Osteomyelitis or septic arthritis from haematogenous spread
    \item Chronic kidney impairment or a prolonged intensive care stay
    \item Recurrence of infection if the source is not controlled
    \item Antifungal side effects and drug interactions, including liver and kidney toxicity
\end{itemize}

\section*{Prognosis and Outcomes}
Candidemia is a serious infection with a high mortality, often reported in the range of 30--40\%. The outlook depends on how quickly treatment starts, the underlying severity of illness, the species of \textit{Candida} and its drug sensitivity, and whether the source of infection can be removed. Among those who survive, recovery may be prolonged, and some people are left with lasting effects from organ failure or from complications such as eye or valve infection. Outcomes are markedly better when antifungal treatment is started promptly and the source is controlled early.

\section*{Prevention and Self-Care}
Prevention is largely a hospital responsibility, since candidemia overwhelmingly affects unwell inpatients:
\begin{itemize}
    \item Meticulous hand hygiene and aseptic technique when inserting and handling intravenous catheters
    \item Removing catheters, lines, and urinary devices as soon as they are no longer needed
    \item Judicious use of broad-spectrum antibiotics to avoid disturbing normal flora
    \item Antifungal prophylaxis in selected high-risk patients, such as certain transplant and chemotherapy patients
    \item Controlling blood sugar in diabetic hospital patients
    \item Using oral or enteral feeding in preference to total parenteral nutrition where possible
\end{itemize}
For people at home, good control of diabetes and general health reduces overall risk. Anyone who has been treated for candidemia should report new fever, red or painful eyes, or chest symptoms to a doctor promptly after discharge.

\section*{Sources and Further Reading}
The following sources provide authoritative, up-to-date information on candidemia and invasive candidiasis.
\begin{itemize}
    \item Centers for Disease Control and Prevention. \textit{Candidiasis and invasive candidiasis.} https://www.cdc.gov/fungal/diseases/candidiasis/
    \item World Health Organization. \textit{WHO fungal priority pathogens list including Candida auris.} https://www.who.int/publications/
    \item MSD Manual. \textit{Invasive candidiasis (candidemia).} https://www.msdmanuals.com/
    \item NHS. \textit{Invasive fungal infections in hospital patients.} https://www.nhs.uk/conditions/
    \item MedlinePlus. \textit{Candidiasis.} https://medlineplus.gov/candidiasis.html
\end{itemize}
